% Source: https://artofproblemsolving.com/community/c3h1779794_3_algebra_questions
% Stable identity: aops-topic-1779794-q03
Each step keeps the total mass $100$ g and removes the fraction $\dfrac{x}{100}$ of the salt present, so it multiplies the concentration by $1-\dfrac{x}{100}$. After two steps
\[
20\left(1-\dfrac{x}{100}\right)^2=5\ \Longrightarrow\ \left(1-\dfrac{x}{100}\right)^2=\dfrac{1}{4}\ \Longrightarrow\ 1-\dfrac{x}{100}=\dfrac{1}{2}
\]
(the value must lie between $0$ and $1$). Hence $x=\boxed{50}$. Check: $20$ g of salt becomes $10$ g, then $5$ g, in $100$ g of water: $5\%$.
